% Model sketch template — macro-model-builder skill
% Macro working-paper defaults, trimmed for a sketch. Sections follow the
% default block map (references/model-blocks.md); delete or merge sections
% freely for small models, but keep: Purpose, Equilibrium, The experiment,
% Open questions.
\documentclass[12pt]{article}

\usepackage[margin=1in]{geometry}
\usepackage{amsmath,amssymb,mathtools,amsthm}
\usepackage{booktabs}
\usepackage{setspace}
\usepackage[round]{natbib}
\usepackage[colorlinks=true,linkcolor=blue,citecolor=blue,urlcolor=blue]{hyperref}

\onehalfspacing

% ---- Theorem-style environments -------------------------------------------
\newtheorem{assumption}{Assumption}
\newtheorem{definition}{Definition}
\newtheorem{proposition}{Proposition}
\newtheorem{lemma}{Lemma}
\newtheorem{corollary}{Corollary}
\theoremstyle{remark}
\newtheorem{remark}{Remark}

% Tag substantive modelling choices, never regularity conditions, as critical
% (mechanism dies without it) or convenience (tractability only):  \begin{assumption}\atype{critical} ...
\newcommand{\atype}[1]{\textnormal{[\textit{#1}]}\;}

% Identifying moment for a parameter, used inline after introducing it:
%   ... with Calvo parameter $\theta$.\moment{$\theta$}{micro price-change frequency}
\newcommand{\moment}[2]{\footnote{Identifying moment for #1: #2.}}

% ---- Metadata --------------------------------------------------------------
\title{Model sketch: MODEL-NAME}
\author{AUTHOR}
\date{\today\ --- draft, not for circulation}

\begin{document}
\maketitle

\section*{Purpose}
% One paragraph, three sentences: the question, the mechanism the model must
% generate, the experiment that will be run. Everything below is checked
% against this paragraph.
[Question.] [Mechanism.] [Experiment.]

\paragraph{Base papers.} \citet{KEY1}; \citet{KEY2}.
% One line each on what is taken from the paper (environment, a trick, a
% calibration) — not a literature review.

\section*{Notation}
% Load-bearing (Thomson): mnemonic letters, no symbol doing double duty,
% defined before use. Delete for small models.
\begin{center}
\begin{tabular}{ll}
\toprule
Symbol & Meaning \\
\midrule
$t$ & time, discrete, infinite horizon \\
% ...
\bottomrule
\end{tabular}
\end{center}

\section{Environment}
% Time, demography, goods, technology at a glance, information, and the
% within-period timeline (write it explicitly if there is any doubt).

\section{Agents and preferences}

\section{Technology and market structure}
% Every price has a known origin: parametric, chosen by an optimizer, or a
% protocol. Every income flow has a recipient.

\section{Agents' problems}
% Problems, FOCs with one-line economic readings, envelope/transversality
% where relevant.

\section{Aggregation and market clearing}
% All clearing conditions; resource constraint derived, checked against the
% income accounting. With heterogeneity: which statistics of the distribution
% enter aggregates.

\section{Equilibrium}
\begin{definition}[Equilibrium]
% The full list of objects and the conditions they jointly satisfy; what is
% exogenous. Mentions everything introduced above and nothing else.
\end{definition}

\section{Steady state}
% Stationary system; which equations pin which variables; the computation
% strategy, down to the one-dimensional root-finding problem if there is one.

\section{The experiment}
% The shock (default: unanticipated at t=0, perfect foresight after), outcome
% variables, and the expected mechanics in words — written before computing.

\subsection*{Parameters and identifying moments}
\begin{center}
\begin{tabular}{lll}
\toprule
Parameter & Role & Identifying moment \\
\midrule
% $\beta$ & discounting & real interest rate in steady state \\
\bottomrule
\end{tabular}
\end{center}

\section{Open questions / next steps}
% Honest and specific: unresolved choices, checks not yet run (pre-flight
% checklist in craft.md), what the quantitative version needs, what the
% sketch demonstrates vs. merely suggests.
\begin{itemize}
  \item
\end{itemize}

% refs.bib beside this file, or wherever BIBINPUTS in latexmkrc points
% (copy assets/latexmkrc beside this file; build with latexmk). AEA style.
\bibliographystyle{aer}
\bibliography{refs}

\end{document}
